% ============================================================
% 分片：s01-intro  对应原文 PDF 第 2--4 页（印刷页码 2--4）
% 负责章节：1. Introduction
% 译者：t01-intro   审校状态：初稿
% 审校：t02-s2a 已完成逐式核对（修正 1 处）
% ============================================================
% 边界：起于第 2 页 "1. INTRODUCTION" 标题；止于第 4 页
% "2. FIELD THEORY BACKGROUND" 标题之前（该标题及 §2 正文归 s02a-fieldtheory）。
% 公式以 pages/page-02.png--page-04.png 为唯一依据逐式核对；第 2--4 页原文无脚注。

% ==== 原文 PDF 第 2 页（印刷页 2）====
\label{op:2}

\section{引言}

\begin{keytakeaway}
\textbf{核心物理图像：}四维 $N=1$ 超对称规范理论中由 Fayet--Iliopoulos $D$ 项引出的参数 $r$，在维度约化到二维之后仍可自由调节；随着 $r$ 改变，同一族二维 $N=2$ 理论在不同的相之间发生相变，超对称可以保持不破缺，但无质量谱发生改变。Calabi--Yau 非线性 $\sigma$ 模型与 Landau--Ginzburg 模型正是同一线性 $\sigma$ 模型的两个相，因此二者的对应可由相结构导出，而不必预先假设共形不变性。

\par\textbf{数学工具：}$D$ 项势能用 $U(1)$ 在 $\mathbb{C}^k$ 上的 Hamilton 函数（矩映射）来刻画，$r$ 是该矩映射的常数平移，真空模空间由相应的辛商给出；此外还用到线性 $\sigma$ 模型、\textit{cc} 手征环、镜像对称与双有理等价类。

\par\textbf{主要结论：}由同一二维规范理论的相结构可以重新得到 C-Y/L-G 对应的几乎全部具体证据，并将其推广到射影空间乘积、环面簇等更一般的 Calabi--Yau 流形以及 $(0,2)$ 世界面超对称模型；FI 参数的变动可以引发保持双有理等价类的 Calabi--Yau 拓扑改变，这为弦论中伴随时空拓扑改变的物理过程提供了首批具体证据。
\end{keytakeaway}

本文的核心是系统地利用一个简单的想法；该想法在四维 $N=1$ 超对称理论中已为人所熟悉，因此下面先在这一语境下陈述它。考虑由矢量（规范）多重态与带荷手征多重态（chiral multiplet）构造的可重整规范场论（gauge theory）。若规范群含有一个 $U(1)$ 因子，则可能的超对称相互作用之一是 Fayet--Iliopoulos $D$ 项。用矢量超场 $V$ 表示，该项可写为

% [Tutor Note: 推导补足 —— (1.1) 与 (1.2)(1.3) 是同一 $D$ 项的两种写法。在超场语言中，$\int \mathrm{d}^4\theta\,V$ 取出矢量超场的 $D$ 辅助场分量，故 (1.1) 给出 $D$ 的线性项 $-rD$；它与规范动能项中的 $\frac{1}{2e^2}D^2$、以及物质场的 $D\sum_i n_i|s_i|^2$ 耦合一起构成作用量中依赖 $D$ 的部分，对 $D$ 变分即得 (1.3)。将 (1.3) 代入 (1.2) 消去 $D$，得 $U=\frac{e^2}{2}\bigl(\sum_i n_i|s_i|^2-r\bigr)^2+\sum_i\bigl|\partial W/\partial s_i\bigr|^2$，即通常的 $D$ 项势加 $F$ 项势；这一步被原文的叙述顺序跳过。]

\begin{equation}\tag{1.1}
  -r\int \mathrm{d}^4x\, \mathrm{d}^4\theta\, V.
\end{equation}

手征多重态 $S_i$ 带荷 $n_i$，其标量分量 $s_i$（$i=1\ldots k$）的势能则为

\begin{equation}\tag{1.2}
  U(s_i)=\frac{1}{2e^2}D^2+\sum_i\left|\frac{\partial W}{\partial s_i}\right|^2,
\end{equation}

其中 $e$ 是规范耦合，$W$ 是全纯函数，称为超势（superpotential），而

\begin{equation}\tag{1.3}
  D=-e^2\left(\sum_i n_i|s_i|^2-r\right).
\end{equation}

实际上，$D$ 可以解释为 $U(1)$ 作用在 $\mathbb{C}^k$ 的一个副本上的 Hamilton 函数（$s_i$ 是该副本上的坐标，其上赋予的凯勒形式为 $\omega=-i\sum_j \mathrm{d}s_j\wedge \mathrm{d}\bar s_j$）。参数 $r$ 对应于给 Hamilton 函数加上一个常数这一熟悉的可能做法。

% [Tutor Note: 几何图像与易错点 —— 设 $U(1)$ 以电荷 $n_i$ 作用于 $s_i$（$s_i\to e^{in_i\alpha}s_i$），则关于 $\omega=-i\sum_j \mathrm{d}s_j\wedge \mathrm{d}\bar s_j$ 的矩映射为 $\mu=\sum_i n_i|s_i|^2$，而 (1.3) 即 $D=e^2(r-\mu)$，故 $D$ 与 Hamilton 函数 $H=r-\mu$ 只相差因子 $e^2$；$r$ 只是把 Hamilton 函数（矩映射）平移一个常数，并不改变 $U(1)$ 的生成元。$D$ 项势 $\frac{1}{2e^2}D^2=\frac{e^2}{2}H^2$ 的零点条件为 $\mu=r$，把该零点集按 $U(1)$ 取辛商即得 $D$ 平坦的模空间。]

% [Tutor Note: 物理图像 —— $D$ 项条件 $\mu=r$ 的解集随 $r$ 的取值而变：在一侧可以在 $s_i\neq0$ 处满足 $D=0$，$U(1)$ 被 Higgs 化（规范场获得质量）；在另一侧（视电荷 $n_i$ 的符号）只能由另外一些场凝聚来满足，无质量谱随之不同。这正是“随 $r$ 改变而发生相变、超对称不破缺而无质量谱改变”在势能层面的机制；在二维中 $r$ 还将与瞬子 $\theta$ 角合成复化凯勒参数，其变动对应靶几何的改变（见 §2、§3）。]

随着 $r$ 改变，这样的系统通常会经历相变（phase transition）；在许多情形下超对称保持不破缺，但无质量场的谱结构发生了改变。这是四维可重整超对称理论唯象学中一个熟悉而重要的故事。本文的目的是研究同样的模型，但把它们维度约化到 $D=2$。这一故事在二维的版本尤为丰富，因为更温和的反常%
% ==== 原文 PDF 第 3 页（印刷页 3）====
\label{op:3}
% [跨页接续: 原文第 2 页末 "because milder anomaly" 与第 3 页首 "cancellation conditions…" 属同一句；本译按语义完整保留，页码锚点置于句中断点。]
抵消条件为构造模型提供了更大的自由度，也因为二维无质量场具有通常的特殊性质。

% [Tutor Note: 推导补足 —— 四维 $N=1$ 超对称经维度约化给出二维 $N=2$（即 $(2,2)$）超对称：四维矢量多重态约化为二维矢量多重态，后者除二维规范场 $A_a$ 外还含有两个实标量（合成一个复标量 $\sigma$）；四维手征多重态约化为二维手征多重态；全纯超势 $W$ 原样保留；四维的 FI 参数 $r$ 直接降为二维的 FI 参数。$\sigma$ 的出现使二维规范理论具有库仑分支与扭超势，这些公式在 §2 中系统给出。]

% [Tutor Note: 物理图像 —— 二维手征规范反常由一圈（两个规范玻色子）图给出，其系数只含电荷的二次组合（左右手费米子的贡献抵消），不像四维那样还须满足 $\sum_i q_i^3=0$ 等条件，故允许的物质内容更丰富。所谓“二维无质量场的特殊性质”，通常指二维无质量标量场没有长程序（Mermin--Wagner--Coleman 定理），其真空期望值可以连续调节，因此相结构比四维更丰富。]

作为一个最初促成本项研究的特例，本文将重新得到 Calabi--Yau 模型与 Landau--Ginzburg 模型之间熟知的对应关系 [1,2,3]。关于这一对应的通常论证包括：普适性（universality）的启发式论证 [1--3]，其中包括 Calabi--Yau 模型与 Landau--Ginzburg 模型均为由相同数据参数化的 $N=2$ 模型这一简单事实；启发式的路径积分论证 [3]；以及 \textit{cc} 手征环（\textit{cc} chiral ring）的等价性——后者在前面刚引用的论文以及 [4-6] 中已有讨论。由本文的构造，可以重新获得关于 Calabi--Yau/Landau--Ginzburg 对应的几乎全部具体证据，并同时得到新的特征。本文的论证既不假设、也不证明 Calabi--Yau 模型或 Landau--Ginzburg 模型的共形不变性（conformal invariance）；若以（众所周知的）关于共形不变性的独立论证来补充，这一图像无疑会更加丰富。

% [Tutor Note: 推导补足与物理图像 —— \textit{cc} 手征环即 (2,2) 超共形场论中手征原初场按 $U(1)_L\times U(1)_R$ 荷分出的 (c,c) 环；其 OPE 系数受超对称保护、在 RG 流下不变，故可在 Calabi--Yau 与 Landau--Ginzburg 两种描述之间直接比较，这是原文把“手征环等价”列为对应证据的原因。同理，本文的论证在紫外（线性 $\sigma$ 模型，并非共形场论）层面比较受保护的算符与相的谱，不依赖任一模型的共形不变性；只有再补充独立的共形不变性论证，比较才提升为红外共形场论层面的等价。]

本文还将找到通常 C-Y/L-G 对应的许多推广，涉及更一般的 Calabi--Yau 流形类（射影空间乘积中的超曲面，更一般的环面簇与其他簇中的超曲面，以及这类超曲面的交）。一个新颖之处是：一般而言，除了普通的 Landau--Ginzburg 模型外，还必须考虑规范 Landau--Ginzburg 模型（gauged Landau--Ginzburg model）。此外，还将把 C-Y/L-G 对应推广到具有 $(0,2)$ 世界面超对称的模型；迄今只考虑过 $(2,2)$ 超对称的情形。$(0,2)$ 世界面超对称具有唯象学意义，因为它可以自然地导致例如具有时空大统一规范群 $SU(5)$ 或 $SO(10)$ 而非 $E_6$ 的模型。

% [背景补充: 环面簇（toric variety）是含代数环面 $(\mathbb{C}^*)^n$ 作为稠密开子集的代数簇，可由扇（fan）的组合数据完全描述；在环面簇中构造超曲面以及镜像流形都可化为多胞形上的组合问题（详见 §5）。]

Fayet--Iliopoulos $D$ 项系数的变动还可以引发涉及 Calabi--Yau 流形拓扑改变（change of topology）的相变。这一点将仅在 §5.5 中从局部观点加以研究。由于 §4 中说明的原因，所讨论的拓扑改变保持 Calabi--Yau 流形的双有理等价类（birational equivalence class）。正是这一情形已由%
% ==== 原文 PDF 第 4 页（印刷页 4）====
\label{op:4}
% [跨页接续: 原文第 3 页末 "Precisely this situation has been studied globally by" 与第 4 页首 "Aspinwall, Greene, and Morrison [7]…" 属同一句；本译按语义完整保留，页码锚点置于句中断点。]
Aspinwall、Greene 与 Morrison [7] 从整体上加以研究：他们借助镜像对称（mirror symmetry）——同时找到拓扑不同但双有理等价、且具有相同镜像的流形对——获得了非常精确而详尽的证据，表明可以存在伴随拓扑改变的物理过程。

% [Tutor Note: 物理图像 —— 对 Calabi--Yau 流形 $X$ 与其镜像 $X^\circ$，相应的 (2,2) 共形场论相同；在此对应下，$X$ 的复化凯勒模（含本文的 $r$）与 $X^\circ$ 的复结构模互换。Aspinwall、Greene 与 Morrison 正是利用这一互换，找到拓扑不同但双有理等价、且共享同一镜像的流形对；由于二者给出同一个 (2,2) 理论，它们对应同一物理理论，这为物理过程可以改变时空拓扑提供了具体证据。]

本文所处理的问题与他们的研究有着出乎意料的密切关系，作者从与他们的讨论中受益，尤其是在 §5.5 方面。[7] 的结果与 §5.5 的结果一起，构成了弦论中（久已猜想的）伴随时空拓扑改变的过程确实会发生的首批具体证据。

本文基于相关 $\sigma$ 模型的熟知表示 [8]，只是所考虑的是线性 $\sigma$ 模型（linear sigma model），而非它们在低能下约化成的非线性 $\sigma$ 模型；这提供了一种有用的额外回旋余地。

% [Tutor Note: 推导补足 —— 线性 $\sigma$ 模型是二维可重整、渐近自由的规范理论，只有在红外极限才约化为 Calabi--Yau 上的非线性 $\sigma$ 模型；保留线性的紫外描述，就不必一开始就写出红外非线性 $\sigma$ 模型（其靶空间描述在相变点会退化），从而获得在 $r$ 上连续变形、穿越相边界的操作余地。这是本文方法的要点。]

在本文工作进行期间，作者收到 Thaddeus 的论文 [9]；该文运用一个与本文密切相关的数学思想，证明了 $SU(2)$ WZW 模型共形块空间维数的 Verlinde 公式。
